\documentclass{article}
\usepackage{amsmath}
\begin{document}

\section*{Machine cycle times: from linear to hierarchical models}

Let $Y_{ij}$ be the cycle time (seconds) of production run $i$ on machine $j$
($j=1,\dots,24$), $X_{ij}$ the hours since the machine's last maintenance, and
$N_{ij}=1$ if the run was on the night shift (0 for day). Machines are spread
over three plants and have between 3 and 20 runs each.

\subsection*{Synthetic data}
The data were simulated as
\begin{align*}
Y_{ij} &= 60 + b_{0j} + (0.05 + b_{1j})\,X_{ij} + 1.5\,N_{ij} + \varepsilon_{ij}\\
b_{0j} &\sim N(0,\ 4^2), \qquad b_{1j} \sim N(0,\ 0.02^2), \qquad
\varepsilon_{ij} \sim N(0,\ 2^2)
\end{align*}
So every machine has its own baseline and wear rate, the night penalty is
the same 1.5\,s for all machines, and there is no true plant effect.

\subsection*{M1: Simple linear model (complete pooling)}
\begin{align*}
Y_{ij} \mid \beta_0,\beta_1,\sigma &\sim N(\beta_0 + \beta_1 X_{ij},\ \sigma^2)
\end{align*}
One line for every machine; only tool wear is used.

\subsection*{M2: Multiple linear model (still pooled)}
\begin{align*}
Y_{ij} \mid \beta,\sigma &\sim N(\beta_0 + \beta_1 X_{ij} + \beta_2 N_{ij}
  + \beta_3\,\text{PlantB}_j + \beta_4\,\text{PlantC}_j,\ \sigma^2)
\end{align*}
Adds shift and plant, but machines within a plant are still treated as identical.

\subsection*{M3: Categorical machine (no pooling of intercepts)}
\begin{align*}
Y_{ij} \mid \beta_{0j},\beta_1,\beta_2,\sigma &\sim
  N(\beta_{0j} + \beta_1 X_{ij} + \beta_2 N_{ij},\ \sigma^2)\\
\beta_{0j} &\sim N(\cdot,\cdot) \ \text{independently for each } j
\end{align*}
Each machine gets its own baseline (23 dummy variables); machines share nothing.

\subsection*{M4: Hierarchical, varying intercepts (partial pooling)}
\begin{align*}
Y_{ij} \mid \beta_{0j},\beta_1,\beta_2,\sigma &\sim
  N(\beta_{0j} + \beta_1 X_{ij} + \beta_2 N_{ij},\ \sigma^2)
  && \text{(within machine)}\\
\beta_{0j} \mid \beta_0,\sigma_0 &\sim N(\beta_0,\ \sigma_0^2)
  && \text{(between machines)}\\
\beta_0,\ \beta_1,\ \beta_2 \sim N(\cdot,\cdot), \quad \sigma &\sim \text{Exp}(\cdot), \quad
\sigma_0 \sim \text{Exp}(1)
\end{align*}
Machine baselines come from a common population, so each is pulled
(shrunk) toward $\beta_0$; machines with few runs are pulled the most.

\subsection*{M5: Hierarchical, varying shift effect}
\begin{align*}
Y_{ij} \mid \beta_{0j},\beta_1,\beta_{2j},\sigma &\sim
  N(\beta_{0j} + \beta_1 X_{ij} + \beta_{2j} N_{ij},\ \sigma^2)\\
\begin{pmatrix}\beta_{0j}\\ \beta_{2j}\end{pmatrix} &\sim
  N\!\left(\begin{pmatrix}\beta_0\\ \beta_2\end{pmatrix},\ \Sigma\right)
\end{align*}
Each machine also gets its own night-shift penalty. In the R code
\texttt{(shift | machine)} has no fixed \texttt{shift} term, so there
$\beta_2 = 0$ and the night penalties are centred on zero.

\subsection*{Model comparison (leave-one-out ELPD; higher is better)}
\begin{tabular}{lrr}
\hline
Model & ELPD diff & SE diff \\
\hline
M4 Hierarchical, varying intercepts & 0.0 & 0.0 \\
M3 Categorical machine & $-0.9$ & 1.4 \\
Linear: hours + shift & $-130.3$ & 12.0 \\
\hline
M4 Hierarchical, varying intercepts & 0.0 & 0.0 \\
M5 Hierarchical, varying shift & $-2.7$ & 4.9 \\
\hline
\end{tabular}

\subsection*{Reading the results}
\begin{itemize}
  \item The pooled linear model ignores the 4\,s spread between machines and
        loses decisively ($-130.3$, about 10 standard errors).
  \item M3 and M4 both capture machine differences and predict about equally
        well. M4 does it with one variance parameter, $\sigma_0$, instead of
        23 dummies, and can predict a \emph{new} machine; M3 cannot.
  \item M5 adds machine-specific night penalties that do not exist in the
        data, so the extra parameters fit noise. Shrinkage keeps the cost
        small ($-2.7$, within one standard error), and the simpler M4 is
        preferred.
\end{itemize}

\end{document}
